
%=======================================================================================

\chapter{Methodology and Experimental Design}
\label{ch:methodology}

This chapter outlines how the raw data was utilised to answer the three research questions. It is organised by method, first the experimental framework, then common design decisions and data leakage control, the external environmental feature set construction, the top-height prediction models, the two attribution analyses, evaluation design, and finally assumptions, limitations and reproducibility.



%------------------------------------------------
\textbf{RQ1} tests which models accurately predict top height and under what input and cohort year. \textbf{RQ2} is an attribution task, asking which areas over or under predict top height relative to the average CR equation formed from the entire dataset. \textbf{RQ3} fits each plot with its own asymptote from only the plot's own repeated observations, and makes comparisons on the growth trajectories of curves, and if they are associated with specific stand, management or external environmental variables.

\begin{table}[!htbp]
\scriptsize
\renewcommand{\arraystretch}{1.2}
\begin{tabular}{p{0.8cm}p{1.4cm}p{3.3cm}p{4.2cm}p{2.9cm}}
\toprule
RQ & Unit & Target & Model families compared & Evaluation design \\
\midrule
RQ1 & Plot-survey row & 95th-percentile LiDAR top height & Chapman-Richards curve, average-by-age
lookup, linear/random forest/XGBoost, deep neural network, Chapman-Richards-guided neural network
variants & Spatial-block (primary), pooled spatial $k$-fold, temporal, plot-level \\
RQ2 & Plot & Mean residual from the single pooled Chapman--Richards curve, across a plot's survey
history & Elastic Net, XGBoost, two-stage mixed-effects model & Spatial-block \\
RQ3 & Plot & Difference between a plot's own fitted asymptote and its yield-class-implied
asymptote & Elastic Net, XGBoost, GNNWR & Spatial-block \\
\bottomrule
\end{tabular}
\caption{Research questions mapped to unit of analysis, target, model families, and evaluation
design.}
\label{tab:rq-framework}
\end{table}

% FIGURE CONTENT (pipeline overview): a flowchart, cohorts feeding data-splitting
% (Table~\ref{tab:splits-summary}) and, separately, the shared Set1-5 screening process
% (Table~\ref{tab:feature-sets-composition-rq1}); both outputs feed three parallel branches, one
% per research question (RQ1 top-height prediction, RQ2 shared-curve residual attribution, RQ3
% plot-specific asymptote attribution); each branch terminates in its own confirmed evaluation
% output (see caption). Do not draw a rendered coefficient/residual map as an output box -- not
% confirmed to exist in the modelling pipeline as of this pass.
\begin{figure}[!htbp]
\centering
\fbox{\parbox{0.9\textwidth}{\centering\vspace{4cm}FIGURE PLACEHOLDER --- pipeline overview, see
comment above source for exact content\vspace{4cm}}}
\caption{Overview of the full pipeline from raw data to each research question's results. Both
cohorts (\S\ref{sec:data-cohorts}) feed the four data-splitting strategies
(Table~\ref{tab:splits-summary}) and a shared screening process that produces each research
question's own Set1--4 (Table~\ref{tab:feature-sets-composition-rq1}). Three branches follow, one
per research question (Table~\ref{tab:rq-framework}): RQ1's top-height prediction models, RQ2's
shared-curve residual attribution, and RQ3's plot-specific asymptote attribution. RQ1 produces
per-fold and pooled predictive metrics under each split type; RQ2 produces a fixed-effects
coefficient table plus held-out residual predictions; RQ3 produces per-fold predictive metrics and
a per-plot coefficient table (GNNWR), together with a residual spatial-autocorrelation diagnostic
(\S\ref{sec:eval-diagnostics}). No branch currently produces a rendered coefficient or residual
map.}
\label{fig:pipeline-overview}
\end{figure}



%------------------------------------------------
\section{Data splits and Feature Design} 
4survey is the primary cohort throughout; 6survey's longer per-plot trajectories suit the temporal split, and its spatial-block results are compared against 4survey's throughout to help explain differences in model behaviour. For RQ1, each row in the dataset represents a single plot in a given survey year. For RQ2 and RQ3, data are aggregated so each row represents a single plot (Table~\ref{tab:rq-framework}). Model performance is evaluated using four distinct data-splitting strategies.

As a held-out plot's nearest neighbour is typically a training plot metres away, a model can succeed by learning \emph{where} it is in Aberfoyle rather than \emph{why} a plot grows the way it does. Spatial block splitting removes this by holding out whole compartments instead. The temporal split asks whether a model trained on past growth still predicts a real future survey it never saw, rather than an unseen location in the same time period. Table~\ref{tab:splits-summary} gives the exact construction of all four splits used across RQ1--RQ3 (Table~\ref{tab:rq-framework}).

\begin{table}[!htbp]
\centering
\scriptsize
\renewcommand{\arraystretch}{1.2}
\begin{tabular}{p{2.1cm}p{5.1cm}p{5.3cm}}
\toprule
Split & Construction & Role \\
\midrule
Plot-level & Individual plots, 60/20/20, random, seed 42 \ct{source that uses plot-level split}. & Baseline reference only, not a leakage-controlled split. \\
Spatial block (primary) & Whole compartments, 60/20/20, row-count-balanced, seed 42; training plots within 60\,m of a held-out plot excluded (val/test always retained) \cite{roberts2017Cross_C5_EVAL}. & Main leakage control for every RQ (Table~\ref{tab:rq-framework}). \\
Spatial block, pooled $k$-fold & 5 compartment folds, row-count-balanced; each held out once as test, next fold as validation, remainder as training; held-out predictions pooled before scoring \cite{ploton2020Spatial_C5_EVAL}. & Stronger confirmation of the single spatial-block split -- every compartment scored exactly once. \\
Temporal (secondary) & Whole survey years assigned to train/val/test: 4survey trains 2008+2012, validates 2021, tests 2023; 6survey trains 2002--2012, validates 2021, tests 2023 \cite{lynch2025Digital_C4_NN}. & Tests generalisation across time (RQ1 only; RQ2/RQ3 have no temporal axis, \S\ref{sec:eval-diagnostics}). \\
\bottomrule
\end{tabular}
\caption{The four data-splitting strategies. All use seed 42; compartment/sub-compartment
identifier columns are never included in any model's feature list, enforced automatically.}
\label{tab:splits-summary}
\end{table}

Every model uses the same split assignment for a given (cohort, split type, seed) combination. Shared partitions and matched architectures (e.g.\ DNN vs.\ CR-PINN) reduce confounding when comparing model families under the same split, though residual differences can still reflect optimisation behaviour or stochastic variation rather than model design alone.

% FIGURE CONTENT (splits figure): three panels left to right on real Aberfoyle geography (reuse
% whatever map base layer the data chapter's Figure~\ref{fig:data-cohort-map} uses, for visual
% consistency). Panel A (plot-level): a small cluster of plots coloured by train/val/test
% membership, with at least one held-out plot's nearest training-plot neighbour explicitly
% connected by a short line/arrow labelled with its distance in metres. Panel B (spatial-block):
% compartment polygons coloured by train/val/test, with the 60m exclusion buffer shown as a
% hatched ring around held-out compartments, and a callout noting buffer-excluded training plots
% are dropped, not reassigned. Panel C: a horizontal timeline bar per cohort (4survey: 2008, 2012
% train -- 2021 val -- 2023 test; 6survey: 2002/2006/2008/2012 train -- 2021 val -- 2023 test),
% with the 2012-2021 gap visually marked.
\begin{figure}[!htbp]
\centering
\begin{minipage}[b]{0.32\textwidth}
  \centering
  \fbox{\parbox{\textwidth}{\centering\vspace{2.5cm}\textbf{Panel A}\\Plot-Level Split\vspace{2.5cm}}}
\end{minipage}\hfill
\begin{minipage}[b]{0.32\textwidth}
  \centering
  \fbox{\parbox{\textwidth}{\centering\vspace{2.5cm}\textbf{Panel B}\\Spatial-Block Split ($60\text{~m}$ Buffer)\vspace{2.5cm}}}
\end{minipage}\hfill
\begin{minipage}[b]{0.32\textwidth}
  \centering
  \fbox{\parbox{\textwidth}{\centering\vspace{2.5cm}\textbf{Panel C}\\Temporal Split Timeline\vspace{2.5cm}}}
\end{minipage}
\caption{Overview of data-splitting strategies across cohorts: (A) Unconstrained random plot-level
split, showing the nearest-training-neighbour leakage mechanism; entire compartments, not
individual plots, are held out in (B) to control spatial leakage, with a $60\text{~m}$ boundary
exclusion buffer that removes training-side plots only (val/test membership is never altered); (C)
sequential temporal split timeline across survey years, with the exact train/val/test years per
cohort shown.}
\label{fig:splits-combined}
\end{figure}

\paragraph*{Preprocessing and leakage prevention}. Continuous features are standardised (zero mean, unit variance) using parameters fitted only on the training partition; age is scaled separately so its training-split standard deviation can be extracted for the physics-loss unit conversion. Categorical soil variables are one-hot encoded. Any plot missing a required feature in any survey year is removed entirely, so no model ever sees an incomplete temporal trajectory -- 39 plots (156 rows) for the terrain-conditioned RQ1 models on 4survey, none on 6survey (affected plot IDs and map location: Appendix~\ref{app:missing-cells}).


%------------------------------------------------
\paragraph*{Environmental feature-set construction}
\label{sec:feature-final}
Candidate environmental and management variables are screened separately for RQ1 (predictive tiers), RQ2 (shared-curve attribution), and RQ3 (plot-level attribution). All three follow the same recipe below, differing only in candidate pool and
target column. This produces four tiered feature sets (Set1–Set4) to evaluate whether expanding environmental inputs improves model performance or introduces unhelpful noise across different model architectures.

\begin{table}[!htbp]
\centering
\scriptsize
\renewcommand{\arraystretch}{1.2}
\begin{tabular}{p{2.6cm}p{9.3cm}}
\toprule
Step & What it does \\
\midrule
1. Coverage screening & Plots missing any candidate variable are dropped entirely: 430 plots from
4survey's candidate pool (413 missing SoilGrids, 39 missing CEH Topographic Wetness Index, mostly
clustered in one compartment; plot IDs in Appendix~\ref{app:missing-env}). 6survey has no missing
values. \\
2. Deduplication & Deterministic duplicates (closed-form functions of another candidate) are
dropped; near-exact empirical duplicates (Spearman $|\rho|\geq0.95$) are reduced to one variable,
preferring externally measured over derived variables. \\
3. Importance ranking & Univariate Spearman $|\rho|$, permutation importance, and drop-column
ablation ($\Delta R^2$ when removed) are combined by rank-averaging, balancing each method's
individual weakness (correlation ignores context; permutation is biased under correlated
predictors; ablation depends on one model/split). \\
4. VIF screening & Variables are iteratively dropped if VIF $>5.0$
\cite{obrien2007Caution_C5_EVAL}; baseline columns are exempt from removal but still inform other
variables' VIF. \\
\bottomrule
\end{tabular}
\caption{The four-step screening protocol applied to build each research question's Set1--Set4
(Table~\ref{tab:feature-sets-composition-rq1}).}
\label{tab:feature-screening-steps}
\end{table}

\begin{table}[!htbp]
\centering
\scriptsize
\renewcommand{\arraystretch}{1.2}
\begin{tabular}{p{1cm}p{5.3cm}p{7.2cm}}
\toprule
Set & Selection rule & RQ1 membership (worked example) \\
\midrule
Set1 & Baseline only. & 4 columns: CanopyCover, time\_since\_thinning,
time\_since\_thinning\_missing, recent\_thinning\_5yr. \\
Set2 & Baseline + the top 10 candidates by combined importance rank (Spearman + permutation
importance + drop-column ablation, averaged); a candidate that fails VIF screening is skipped and
replaced by the next-ranked candidate, so the count stays at 10. & 14 columns: Set1 + (e.g.\
gwa\_weibull\_k\_50m, dist\_to\_scpt\_boundary; full list Appendix~\ref{app:feature-sets}). \\
Set3 & Baseline + the upper half (by combined rank) of the terrain/wind category only -- ranked
within that category, not against every other candidate; VIF-screened. & 15 columns: Set1 + (e.g.\
gwa\_weibull\_k\_50m, elevation; full list Appendix~\ref{app:feature-sets}). \\
Set4 & Set3 + the upper half (by combined rank) of every other category -- climate, soil/site,
spatial-position/edge-effects; VIF-screened. & 18 columns: Set3 + (e.g.\ dist\_to\_scpt\_boundary,
tas\_mean; full list Appendix~\ref{app:feature-sets}). \\
\bottomrule
\end{tabular}
\caption{The four feature sets generated by the screening protocol, using RQ1 as an example. While RQ2 and RQ3 follow the identical protocol, differing candidate pools and target variable for the models lead to different set membership (full lists in Appendix~\ref{app:feature-sets}). Set4 is deliberately the final, richest tier: it is the widest set still curated by evidence-driven relevance ranking, extending Set3's terrain/wind-only ranking to every candidate category.}
\label{tab:feature-sets-composition-rq1}
\end{table} %Layout of table not good and third coloumn not filled in with rq1 worked example. 

%------------------------------------------------
\section{RQ1: Top-height prediction models}
\label{sec:rq1-models}

RQ1 evaluates top-height models across four dimensions: baseline predictive accuracy, the use of supplementary environmental data (which can help stakeholders learn more where to invest into further high-quality data acquisition to improve modelling and planning), the benefit of Chapman-Richards physics guidance over unguided networks, and generalisation across spatial and temporal splits.

\textbf{Chapman-Richards (CR) curve.} A three-parameter curve
\begin{equation}
H(a) = y_{\max}\left(1-e^{-ka}\right)^{p}
\label{eq:cr}
\end{equation}
predicts top height $H(a)$ at age $a$ using asymptotic height $y_{\max}$, growth rate $k$, and shape parameter $p$. Curves are fitted strictly on each split's training partition via nonlinear least squares (five restarts, lowest sum of squared errors (SSE); Appendix~\ref{app:sse}). Under the pooled $k$-fold split, five separate curves are fitted per cohort (one per fold). The resulting $(y_{\max}, k, p)$ serve as fixed constants in the PINN loss: each PINN model uses strictly its own fold-specific curve so that information from held-out validation or test compartments never leaks into the physics guidance loss. The CR form assumes a single monotonic growth trajectory per plot; a plot whose real trajectory is non-monotonic (e.g.\ following disturbance) is not represented by this functional form, and is instead handled by the outlier audit described later in this chapter rather than by the curve itself. %how does this related to forwards and backwards pass?

Three standard supervised models are fit on all RQ1 feature sets (baseline and Sets 1–4) using identical splits.
\textbf{Linear regression.} A single linear relationship was fitted between the predictors and height. Because plots were at least 15 years old in 2008, the observed section of the age--height relationship may appear approximately linear despite the nonlinear full growth trajectory. Linear regression provides a simple baseline but performs poorly near the limits of the observed age range (Figure~\ref{fig:age-linearity}).

% FIGURE CONTENT (age/CR-linearity): training-partition age distribution as a histogram (both
% cohorts, or 4survey only if space-limited), with the pre-2008 age floor and the associated
% "unreliable LiDAR scan" exclusion band annotated directly on the x-axis. Overlay a single fitted
% Chapman-Richards curve (Equation~\ref{eq:cr}) across the same age range, with the observed age
% window (roughly 15+ years, per the age-floor stated in the data chapter) highlighted to
% show visually that it sits on the curve's near-linear early-to-mid segment.
\begin{figure}[!htbp]
\centering
\fbox{\parbox{0.9\textwidth}{\centering\vspace{4cm}FIGURE PLACEHOLDER --- training-partition age
distribution with fitted CR curve overlay, see comment above source for exact
content\vspace{4cm}}}
\caption{Training-partition age distribution against the fitted Chapman--Richards curve
(Equation~\ref{eq:cr}). The observed age range covers only part of the full sigmoid trajectory,
sitting on its near-linear early-to-mid segment -- the reason linear regression is a defensible
baseline despite the true age--height relationship being nonlinear.}
\label{fig:age-linearity}
\end{figure}

\textbf{Random Forest.} An ensemble of bootstrap-aggregated decision trees that models non-linear interactions without imposing parametric growth constraints. This introduces a non-parametric inductive bias distinct from both Chapman-Richards curves and DNN/PINN architectures, testing whether top-height prediction requires smooth functional constraints or benefits from flexible ensemble averaging \cite{lee2018Machine_C3_LID}

\textbf{XGBoost} builds regression trees sequentially, with each tree correcting errors remaining from the preceding ensemble:
\begin{equation}
\hat{y}_i^{(K)} = \sum_{k=1}^{K}\eta f_k(\mathbf{x}_i),
\end{equation}
where $\mathbf{x}_i$ contains the predictors for observation $i$, $f_k$ is the $k$th regression tree, and $\eta$ is the learning rate. Unlike Random Forest, which averages independently fitted trees, this sequential approach can capture complex nonlinear relationships and provides a strong conventional machine-learning comparison for the neural models.

To keep hyperparameters comparable across research questions, a single XGBoost configuration (max depth 4, learning rate 0.04, up to 500 boosting rounds with early stopping) is used throughout, carried over from an earlier exploratory fit. RQ3 used this configuration from the outset.  (Appendix~\ref{app:hyperparameters}).


% FIGURE CONTENT (model architecture): a block diagram. Left: [Age, no-env features] -> main trunk
% (3x128, shared) -> [predicted height]. Below/beside: [Terrain/wind features] -> small asymptote
% sub-network (16 units) -> [y_max adjustment, additive] and, dashed/optional for the
% rate-conditioned variant, -> rate sub-network -> [k adjustment, multiplicative, log-space]. A
% separate branch shows [predicted height] and [age] feeding into an autograd box labelled dH/da,
% compared against a box labelled "CR analytical derivative (frozen y_max, k, p)" to produce the
% physics loss; a second small diagram shows two real survey observations of one plot feeding two
% forward passes, differenced and compared to the same CR derivative at mid-age, to produce the
% trajectory loss. These survey pairs are pre-built once from real consecutive-survey rows before
% training starts (not synthesised per batch) -- draw them as a small fixed table/lookup feeding
% the diagram, not as an on-the-fly sampling step.
\begin{figure}[!htbp]
\centering
\fbox{\parbox{0.9\textwidth}{\centering\vspace{4cm}FIGURE PLACEHOLDER --- Model Architecture Block Diagram\vspace{4cm}}}
\caption{Shared architecture and loss formulation for the DNN and CR-PINN's. All variants share a 3-layer predictive trunk ($3 \times 128$). In environment-conditioned models, terrain and wind features feed only auxiliary sub-networks (16 units) to adjust the asymptotic height ($y_{\max}$) and growth rate ($k$). The guided variants augment the standard data loss with: (1) an autograd-derived \emph{physics loss} ($\mathrm{d}H/\mathrm{d}a$) evaluated against the analytical derivative, and (2) a \emph{trajectory loss} computed over consecutive empirical survey pairs.}
\label{fig:model-architecture}
\end{figure}

All DNN and CR-PINN models share an identical core architecture (Figure~\ref{fig:model-architecture}) so that performance differences reflect the loss formulation rather than variations in network capacity \cite{kovacs2022Conditional_C4_NN,miao2023VCPINN_C4_NN}. Environmental covariates enter the DNN directly, whereas in the CR-PINN they are passed through compact auxiliary networks that condition $y_{\max}$, $k$ and $p$ testing whether site and stand variations alter the overall CR curve at the plot level \cite{scolforo2013Dominant_C2_CR}. A network-size sensitivity check across four alternative sizes, both cohorts, tests whether this shared size is well-chosen (Appendix~\ref{app:architecture-screening}). 


The \textbf{Deep Neural Network} uses three hidden layers of 128 units each, leaky-ReLU activation, age concatenated with every other
input feature. The DNN's loss is plain mean squared error,
\begin{equation}
\mathcal{L}_{\text{data}} = \frac{1}{N}\sum_{i=1}^{N}\left(\hat H_i - H_i\right)^2,
\label{eq:data-loss}
\end{equation}
where $\hat H_i$ is the network's predicted height for row $i$ and $H_i$ is the observed height,
plus an L1 weight penalty ($\lambda_{L1}=10^{-5}$). The environment-conditioned DNN is
architecturally identical to the no-environment DNN; its terrain/wind block is simply concatenated
into the same input, widening that vector rather than adding a second network. This is so that interactions between age and stand variables can be learnt without imposing a separate structure.

The \textbf{Chapman-Richards-guided neural network (CR-PINN)} models differ from the baseline DNN only by their loss function. In the baseline no-environment variant, the CR parameters $(y_{\max}, k, p)$ are treated as fixed cohort constants (Equation~\ref{eq:cr}) rather than learned parameters. The model optimises a three-term loss informed by the analytical growth derivative $H'(a) = y_{\max} p k e^{-ka}(1 - e^{-ka})^{p-1}$:

\begin{subequations}
\label{eq:cr-pinn-losses}
\begin{align}
\mathcal{L}_{\text{phys}} &= \frac{1}{N}\sum_{i=1}^{N}\left(\frac{\partial \hat{H}_i}{\partial a_i} - H'(a_i)\right)^2, \label{eq:physics-loss} \\
\mathcal{L}_{\text{traj}} &= \frac{1}{M}\sum_{j=1}^{M}\left(\frac{\hat{H}_{j,\text{later}} - \hat{H}_{j,\text{earlier}}}{\Delta a_j} - H'(\bar{a}_j)\right)^2, \label{eq:trajectory-loss} \\
\mathcal{L} &= \mathcal{L}_{\text{data}} + \lambda_{\text{phys}}\mathcal{L}_{\text{phys}} + \lambda_{\text{traj}}\mathcal{L}_{\text{traj}} + \lambda_{L1}\sum_{\theta}|\theta|. \label{eq:total-loss}
\end{align}
\end{subequations}

Here, the physics loss \eqref{eq:physics-loss} checks the model's instantaneous growth rate via automatic differentiation at single points (rescaled by the training-split standard deviation to match real-age units). The trajectory loss \eqref{eq:trajectory-loss} checks the predicted height gain across $M$ consecutive real survey visits from the same plot, where $\Delta a_j$ is the time elapsed and $\bar{a}_j$ is the mid-point age (drawn strictly from training rows to avoid leakage). The total objective \eqref{eq:total-loss} combines data loss, both consistency penalties (default $\lambda_{\text{phys}}=\lambda_{\text{traj}}=1.0$; ablated below), and $L_1$ regularisation. Because \eqref{eq:physics-loss} and \eqref{eq:trajectory-loss} pull predictions toward the pooled curve's own derivative, this constraint could in principle bias predictions toward a misspecified reference curve rather than the true growth process; the physics-weight ablation tests how much the constraint actually shapes predictions. 

For the \textbf{Environment-conditioned growth} variants, terrain and wind features $\mathbf{x}^{\text{terrain}}_i$ feed into compact sub-networks (16 units) to produce plot-specific parameter adjustments, leaving the shape parameter $p$ fixed as a global constant:

\begin{subequations}
\label{eq:env-adjustments}
\begin{align}
y_{\max}^{(i)} &= y_{\max} + f_{y}\!\left(\mathbf{x}^{\text{terrain}}_i\right), \label{eq:ymax-adjust} \\
k^{(i)} &= k \cdot \exp\!\left(f_{k}\!\left(\mathbf{x}^{\text{terrain}}_i\right)\right). \label{eq:k-adjust}
\end{align}
\end{subequations}

Here, $f_y$ applies an additive adjustment to the asymptotic ceiling $y_{\max}$ \eqref{eq:ymax-adjust} following site-index modelling principles \citep{socha2023Higher_C1_SS}, while the rate-conditioned variant adds $f_k$ to apply a multiplicative adjustment to the growth rate $k$ \eqref{eq:k-adjust}. Both sub-networks receive only environmental features and are initialised near zero, ensuring training begins from the validated pooled curve before learning local adjustments.

\paragraph*{Hyperparameter sweeps and ablation.}
To test the impact of physical guidance, penalty weights are varied across $\lambda_{\text{phys}} = \lambda_{\text{traj}} \in \{0, 1, 2\}$ against an identically sized baseline DNN. The default weight ($\lambda = 1.0$) in \eqref{eq:total-loss} was established prior to evaluation rather than tuned post-hoc, avoiding circularity. All headline models maintain this fixed setting, with detailed ablation outcomes provided in \S\ref{sec:results}.
\begin{table}[!htbp]
\scriptsize
\renewcommand{\arraystretch}{1.2}
\begin{tabular}{p{3.2cm}p{3cm}p{2.6cm}p{4.5cm}}
\toprule
Model & Environmental input & Loss terms & Notes \\
\midrule
CR (baseline) & None & N/A (curve fit) & Pooled per cohort, refit per fold for pooled $k$-fold. \\
Linear / random forest / XGBoost & No-environment feature set & Model-specific & XGBoost uses a
fixed configuration, not a validation-selected grid (Appendix~\ref{app:hyperparameters}). \\
DNN (no environment) & None & Data (Eq.~\ref{eq:data-loss}) + L1 & Shared architecture with every
guided variant. \\
DNN (environment-conditioned) & Terrain/wind, concatenated into main input & Data + L1 & Same
network class, wider input only. \\
Guided network (no environment) & None & Data + physics + trajectory + L1 & $y_{\max},k,p$
pooled, frozen. \\
Guided network (terrain-conditioned asymptote) & Terrain/wind, asymptote sub-network only & Data +
physics + trajectory + L1 & Eq.~\ref{eq:ymax-adjust}; $k,p$ pooled. \\
Guided network (terrain-conditioned asymptote and rate) & Terrain/wind, asymptote and rate
sub-networks & Data + physics + trajectory + L1 & Eqs.~\ref{eq:ymax-adjust}--\ref{eq:k-adjust};
$p$ pooled. \\
\bottomrule
\end{tabular}
\caption{Model family comparison for top-height prediction (RQ1). Full hyperparameters:
Appendix~\ref{app:hyperparameters}.}
\label{tab:rq1-models}
\end{table}

%------------------------------------------------
\section{RQ2a: Environmental Correction of Curve Deviation}
RQ2a reuses CR-parameter predictions from RQ1 alongside the fold-matched CR baseline without fitting new models. Deviations from this baseline curve mark plots that diverge from the population average. We test whether environmental covariates reduce these departures, particularly for the most severely over- or under-predicted plots, rather than providing a uniform global improvement.

For each test observation, predictions from the fold-matched CR baseline and the environment-conditioned models (DNN, PINN, and PINN$_k$) are matched by plot and survey year. The residual reduction is computed from the error magnitudes:
\begin{equation}
\text{reduction}_i = \left|r_i^{\text{CR}}\right| - \left|r_i^{\text{model}}\right|,
\label{eq:residual-reduction}
\end{equation}
where a positive value indicates the model predicts closer to the true height than the baseline curve. Evaluating the residual's sign ($r_i = y_i - \hat{y}_i$) alongside \eqref{eq:residual-reduction} identifies whether the model corrects directional over- or under-prediction by the shared curve, rather than merely shrinking error magnitude.

To separate general model improvements from specific environmental corrections, test observations are partitioned within each fold into equal quartiles by baseline error magnitude $|r^{\text{CR}}|$ (from Q1: smallest error, to Q4: largest). Because this split is purely rank-based %what split? why is it rank based?
, each quartile contains exactly 25\% of the test data; this evaluates performance across standard subgroups rather than focusing on a few extreme outliers. Mean reductions are evaluated per quartile across all three environment-conditioned models, feature sets (Sets~2--4), and cohorts under spatial cross-validation, with the top-performing configuration verified across multiple random seeds for robustness %%NOT ACTUALLY SURE I DID THIS, CHECK WITH RESULTS, POSSIBLY CHANGE PARAGRAPH ON HOW WE EVALUATED RQ2a

%------------------------------------------------
\section{RQ2b: Persistent Shared-Curve Deviation}
\label{sec:rq2b}
This asks which environmental/management conditions align with a plot's persistent position above or below the single shared curve, using only the 4survey cohort (6survey's 47 compartments were judged too few for a reliable spatial-block attribution split).

\textbf{Target construction.}
For each survey row $i$ in a plot's history, a residual is computed against the pooled baseline curve, and the attribution target $\overline{r}_{\text{plot}}$ is defined as the plot's mean residual across all $T$ observed surveys:

\begin{subequations}
\label{eq:target-construction}
\begin{align}
r_i &= H_i - H(a_i), \label{eq:cr-residual} &
\overline{r}_{\text{plot}} &= \frac{1}{T}\sum_{t=1}^{T} r_{i,t}. \label{eq:mean-cr-residual}
\end{align}
\end{subequations}

Age is excluded from all candidate feature sets because it is the baseline curve's sole predictor in \eqref{eq:cr-residual}, which would introduce circularity into the target in \eqref{eq:mean-cr-residual}.

Three models are fitted on $\overline{r}_{\text{plot}}$ across feature sets (Sets~2--4) under spatial-block cross-validation. Set~1 (baseline only) serves as a targeted ablation to test whether canopy cover signals reflect genuine environmental drivers or prior growth.

\begin{enumerate*}[label=(\roman*)]
\item \textbf{Elastic Net} standardises continuous features, one-hot encodes soil classes, and tunes regularisation parameters ($\ell_1/\ell_2$ ratio and penalty strength) via 5-fold CV on the training partition.
\item \textbf{XGBoost} a separate model from RQ1's, fitted on RQ2's own target
  $\overline{r}_{\text{plot}}$ and RQ2's own candidate features;
  what carries over is only the hyperparameter recipe, not the fitted model itself (500 boosting
  rounds, maximum tree depth 4, learning rate 0.04, early stopping against a held-out validation fold).
\item \textbf{Two-stage mixed-effects model.} Uses a two-stage approximation to a full non-linear mixed-effects model \citep{socha2016Assessment_C2_CR}. Stage~1 isolates persistent plot-level departures ($\overline{r}_{\text{plot}}$) from the frozen baseline curve, and Stage~2 fits a linear mixed-effects model against candidate environmental and management features:
\begin{equation}
\overline{r}_{\text{plot}} = \boldsymbol{\beta}^{\top}\mathbf{x}_{\text{plot}} + u_{\text{cpmt}} + \varepsilon_{\text{plot}},
\label{eq:mixed-effects}
\end{equation}
where $\boldsymbol{\beta}$ represents fixed-effect coefficients on standardised features $\mathbf{x}_{\text{plot}}$, $u_{\text{cpmt}} \sim \mathcal{N}(0, \sigma_u^2)$ is a compartment-level random intercept accounting for spatial clustering, and $\varepsilon_{\text{plot}}$ is the residual error. The model uses random intercepts only; random slopes were omitted because compartment sizes vary widely (from single plots to over a thousand) and no varying-effect hypotheses were posed. Fitted on the training partition just for attribution rather than held-out prediction, coefficients are estimated via restricted maximum likelihood (REML). Nested fixed-effect comparisons use maximum likelihood (ML), as REML likelihoods are not comparable across differing fixed-effects structures \citep{pinheiro2000Mixed_C5_EVAL}. Because $\overline{r}_{\text{plot}}$ is a pre-calculated residual, uncertainty from Stage~1 is not propagated; coefficients are interpreted strictly as statistical attribution of curve departures rather than causal mechanisms.
\end{enumerate*}

Feature importance and SHAP values provide post-hoc interpretability, compared against the mixed-effects coefficients for cross-model convergence. %todo, CHECK WE ACTUALLY DO THIS IN THE RESULTS. 

%------------------------------------------------
\section{RQ3: Plot-Specific Curve Deviation}
\label{sec:rq3-target}
This tests whether a plot's fitted growth ceiling deviates from its static yield-class expectation, and whether local environmental or management conditions explain that deviation. Deviation is estimated in the asymptotic ceiling alone, with the curve's shape parameters held fixed, for three reasons: with few survey observations per plot, fixing shape lets the asymptote be estimated by a stable closed-form linear fit rather than an unstable nonlinear optimisation; the yield-class comparison concerns maximum height potential directly, making the asymptote the parameter of interest; and this follows standard site-index modelling convention, where site quality shifts asymptotic potential while trajectory shape stays constant \citep{socha2021Regional_C2_CR}. By construction the target captures asymptotic height potential alone -- plots with different growth rates or inflection timing that ultimately reach the same ceiling are not treated as divergent.


For each plot, the asymptote $\hat{y}_{\max}^{\text{plot}}$ is estimated from its own repeated height--age surveys via closed-form weighted least squares through the origin, fixing the curve's shape parameters to their Forestry Commission yield-class lookup values. The target is the deviation from the static yield-class asymptote:
\begin{equation}
\Delta y_{\max}^{\text{plot}} = \hat y_{\max}^{\text{plot}} - y_{\max}^{\text{yldc}}.
\label{eq:local-ymax-diff}
\end{equation}
This metric is constructed strictly from height--age records and static tables, with no environmental variables involved.

Plots with extreme $\Delta y_{\max}^{\text{plot}}$ values are audited before modelling against two independent criteria: disturbance flags (clearfelling, measurement anomalies, ambiguous disturbance; exact rules in Appendix~\ref{app:disturbance-rules}) and a trajectory-instability cutoff (top 1\% within-plot height change). This audit's criteria are fixed before any model is fitted; its outcome -- which plots, if any, it flags, and whether they are excluded -- is an empirical result and is reported with RQ3's findings in the Results chapter, not asserted here. %MAY REMOVE OR DECREASE COMMENT SIZE AS 0 OF THESE WHERE FOUND. 

Elastic Net and XGBoost are then fitted on $\Delta y_{\max}^{\text{plot}}$ using Sets~2--4, following the same protocol as RQ2b \citep{zou2005Regularization_C4_NN, chen2016XGBoost_C4_NN}. To test if the environment-deviation relationship varies across the landscape rather than following one global trend, Geographically and Neurally Weighted Regression (GNNWR) is additionally fitted \citep{du2020Geographically_C6_GWR}:
\begin{equation}
\hat{y}_i = \beta_0(s_i) + \sum_j \beta_j(s_i)\,x_{ij},
\label{eq:gnnwr}
\end{equation}
where $s_i$ is plot $i$'s spatial coordinate and each $\beta_j(s_i)$ is a spatially varying coefficient. GNNWR learns how these coefficients vary spatially: distances from each plot to a fixed training reference set feed a neural network that outputs multiplicative weights on a frozen global OLS fit, giving each plot its own local coefficients $\beta_j(s_i)$ (Figure~\ref{fig:rq3-attribution}). This directly tests spatial non-stationarity in the environment-deviation relationship, which a single global model cannot represent by construction.

XGBoost attribution uses gain importance and SHAP values \citep{lundberg2017}. Unlike RQ2b's full-sample SHAP, RQ3's SHAP values are computed out-of-fold across the five spatial-block splits and pooled, so each plot's attribution comes from a model that never trained on its own fold. 

%MAY NEED TO ADD IN THE OTHER THINGS WE CALCUALTED AND HOW TO SEE IF DEVATIONS WHERE STORMS OR CATEOGRY VARIABLES ETC HERE, ITS NOT CURRENTLY MENTIONED. 


\begin{figure}[!htbp]
\centering
\fbox{\parbox{0.9\textwidth}{\centering\vspace{4cm}FIGURE PLACEHOLDER --- RQ3 target construction
and attribution models\vspace{4cm}}}
\caption{Plot-specific asymptote-deviation target construction (RQ3) and its three attribution
models. Left to right: the closed-form per-plot asymptote fit vs.\ the static yield-class
asymptote (Equation~\ref{eq:local-ymax-diff}) feeds Elastic Net, XGBoost, and GNNWR; GNNWR's box
additionally shows the distance-to-reference-set input, the spatial weighting network, and the
per-predictor weights multiplied against the frozen global OLS coefficients. All three converge on
a spatial-block holdout evaluation; RQ2b's parallel but structurally simpler pipeline (pooled
residual $\to$ Elastic Net/XGBoost/two-stage mixed-effects model) is described in prose and
equations only earlier in this chapter, not duplicated here, since it has no comparable internal
mechanism to unpack. Annotate: association, not causation.}
\label{fig:rq3-attribution}
\end{figure}
% FIGURE CONTENT (verified against the actual GNNWR implementation before drawing):
% - the spatial weighting network's input is a vector of Euclidean distances from the focal plot
%   to every plot in the reference set (MinMax-scaled), NOT raw coordinates or coordinate
%   differences; no k-NN or distance-threshold step exists, "neighbourhood" is learned implicitly
%   from the full distance vector.
% - the network outputs one multiplicative weight per predictor (+ intercept), NOT final
%   coefficients directly; this is multiplied against a separately pre-fitted, frozen global OLS
%   coefficient to give beta_j(s_i). One shared multi-output network produces every coefficient's
%   weight jointly, not one network per coefficient.
% - held-out/test plots reuse the training reference set for their distance vector (not a bespoke
%   neighbour set built from their own location); the reference set is the full, uncapped training
%   partition in the reported runs (confirmed: run with --reference-set-size 0).
% - no geographical kernel or bandwidth parameter exists in the implementation -- do not draw one.
% - the only confirmed outputs are a per-plot coefficient table and a Moran's I residual
%   diagnostic -- label it "per-plot coefficient table", not "coefficient map", unless a
%   map-producing script is separately confirmed.

\section{Model Evaluation and Robustness}
\label{sec:eval-diagnostics}
This section states what is measured and why, not what was found (Results chapter).

\textbf{Predictive diagnostics (RQ1): }
Every RQ1 model is scored on residuals $e_i = H_i - \hat H_i$ over $N$ rows: MAE $=\frac{1}{N}\sum_i|e_i|$, RMSE $=\sqrt{\frac{1}{N}\sum_i e_i^2}$, $R^2 = 1-\sum_i e_i^2/\sum_i(H_i-\bar H)^2$ (variance explained relative to a mean-only baseline; 1 perfect, 0 matches the mean, negative is worse than the mean on held-out data), and bias $=\frac{1}{N}\sum_i e_i$ (positive is systematic under-prediction). Metrics are further broken down by five age bands to check whether an aggregate score conceals age-specific weaknesses (the maturity gate is applied at plot level, so individual rows below 30 can still appear; a poor $<30$ score reflects a small, non-representative sample, not necessarily worse prediction) %TODO: not sure we did this or if this age thing was important. 
Pooled $R^2$ additionally carries a cluster-bootstrap 95\% confidence interval, resampling whole compartments rather than individual rows or plots.

\textbf{Spatial autocorrelation (Moran's $I$ and the semivariogram range):}
Global Moran's $I$ for model residuals $e_i = y_i - \hat y_i$ tests whether nearby plots have more similar residuals than distant ones, against a null of spatial randomness:
\begin{equation}
    I = \frac{N}{W} \frac{\sum_{i=1}^{N} \sum_{j=1}^{N} w_{ij} (e_i - \bar{e})(e_j - \bar{e})}{\sum_{i=1}^{N} (e_i - \bar{e})^2}
    \label{eq:morans_i}
\end{equation}
where $\bar e$ is the mean residual, $w_{ij}$ the spatial weight between plots $i$ and $j$, and $W=\sum_{i,j} w_{ij}$ (weighting scheme: Appendix~\ref{app:morans-weights}). Values near zero indicate spatial randomness, positive values residual clustering, negative values dispersion. Moran's $I$ is used two ways: as a post-fit diagnostic on RQ1's raw-height residuals (purely diagnostic, never fed back into feature selection), and as a category-removal check for RQ3 (does removing an environmental category increase residual clustering, currently computed on Set4/4survey only). It has not yet been computed on RQ2b's or RQ3's main attribution-model residuals. %in background say why this is useful for spatial attibtuion? 
A semivariogram, fitted to test-partition residuals only, complements Moran's $I$ with the distance range over which spatial autocorrelation persists, rather than only whether it is present.

Feature importance for RQ2b and RQ3, is evaluated using XGBoost gain and SHAP values \citep{lundberg2017unified}, which decompose individual predictions into signed, additive feature contributions analogous to linear model coefficients. Neither metric is used for feature selection. %todo, LETS JUST DOUBLE CHEKC, IF XGBOOST IS THAN REMOVE XGBOOST GAIN, WE MAY NOT EVEN USE IT IN RESULTS. 
In RQ2b, SHAP provides a non-linear global attribution view to assess convergence alongside Elastic Net and mixed-effects coefficients; in RQ3, SHAP additionally enables local attribution for individual outlier plots, isolating specific drivers of extreme growth deviations. In both, SHAP values are pooled across the five spatial folds to report mean feature contributions; fold-to-fold SHAP stability is not evaluated.

Elastic Net coefficients are reported as mean $\pm$ standard deviation across the five spatial folds in both RQ2b and RQ3, to confirm parameter signs and magnitudes remain consistent across spatially distinct partitions rather than reflecting one data split. Attribution models are evaluated across spatial folds without multi-seed sweeps.

\textbf{Robustness to stochastic and configuration choices:} %TODO: CHECK if this was done or needs to be updated. 
\begin{itemize}
\setlength\itemsep{0pt}
\item \textbf{Training-seed reseed} (RQ1, RQ2a): the RQ1 winner is refit across 4 additional seeds; RQ2a reuses these predictions rather than reseeding separately.
\item \textbf{Architecture-size sensitivity} (RQ1, all three guided-network variants): tests whether the fixed network size affects the headline comparison.
\item \textbf{Physics-weight ablation} (RQ1, PINN and PINN$_k$): tests whether the guided network's advantage, if any, depends on the physics/trajectory loss weighting rather than the mechanism.
\item \textbf{XGBoost hyperparameter sensitivity} (RQ1 baseline): tests whether RQ1's headline comparison depends on the library-default-vs-shared-configuration choice; reported as a sensitivity check alongside the original fit, not yet adopted into RQ1's production comparison (outcome in the Results chapter). RQ2b's own XGBoost model has already been refit under the shared configuration in production, so it needs no separate sensitivity check here.
\item \textbf{Split-seed reseed} (RQ3, GNNWR, 6survey only): tests whether GNNWR's smaller-cohort result is stable across different fold assignments, discussed as a limitation below.
\end{itemize} 



\section{Assumptions and Limitations}
\label{sec:limitations}

\begin{itemize}
\setlength\itemsep{2pt}
\item The 60\,m  spatial buffer is set from typical plot spacing, not from the residual's own measured spatial-autocorrelation range. A semivariogram-based buffer could tighten or loosen this bound, but was not used to set the buffer itself because the semivariogram is computed from held-out residuals after a split already exists, a circular basis for choosing that same split's buffer. Compartment-level holdout remains the primary leakage control; the buffer is a secondary refinement on top of it.
\item Feature-set membership is fixed once, not re-derived per fold. Set1--4 are ranked and VIF-screened once on the full training partition, not separately per cross-validation fold. A fold-specific ranking would better separate selection variance from model variance, but was not used because it would let each of the five folds train on a different feature set, breaking the direct model-family comparison every research question in this chapter relies on.
\item Repeated observations are not modelled as within-plot correlated data. RQ1's modelling unit is the plot-survey row; the pooled-$R^2$ bootstrap CI resamples whole compartments, capturing between-plot spatial clustering, but does not separately model correlation between a single plot's own repeated surveys. This is retained because compartment-level resampling is the clustering that matters for RQ1's spatial-generalisation question; a full hierarchical plot-within-compartment model was judged unnecessary complexity for a predictive-accuracy comparison.
\item Environmental variables are point-extracted onto plots smaller than their source resolution. Terrain/wind/soil rasters (sources and resolutions: Chapter~\ref{ch:data}) are matched to each plot's centre; several nearby plots can therefore share an identical value from the coarsest sources. This reduces the effective variation available to the importance-ranking and VIF steps for those variables specifically.
\item XGBoost's hyperparameter recipe is shared across research questions with different targets. The same configuration is applied to RQ1's height target, RQ2b's mean curve-residual target, and RQ3's asymptote-deviation target, rather than tuned per target. This keeps the three research questions' XGBoost models directly comparable, but means only RQ3 is individually tuned for its own target's scale or noise structure \fc{check its was tuned only to rq3}.
\item Both attribution designs (RQ2b's fixed-effects regression, RQ3's GNNWR/Elastic Net/XGBoost) report statistical association between environmental/management conditions and curve departure but cannot distinguis if it is causal. 
\end{itemize}

%%I WILL REVIEW THIS AFTER I DRAFT RESULTS AND BACKGROUND, AND I IMAGINE MUCH OF THE SIMPLE METHOD JUSTIFCAITON CAN GO TO BACKGROUND, (SIMPLIFY WHAT IS SAID IN METHODLOGY WHITHOUT EXPLINAING), BUT THEN ADD IN MORE EVLAUATION RULES AND AFTER MODEL RUN DIAGNOSTICS THAT IM STILL DOING. 

%------------------------ rq1_ADDTIONS ------------------------
%% tuned xgboost separetaly so now its no the same as rq3, its for this task (rq2 isnt tuned at the moment)
%% RQ1 XGBoost now has its own tuned config, not the shared RQ2b/RQ3 one -- update:
%%   - Table~\ref{tab:rq1-models} XGBoost row: "fixed configuration, not a validation-selected grid" -> now IS validation-selected, per-cohort
%%   - the "single XGBoost configuration...used throughout" sentence (Sec RQ1, para on XGBoost): no longer true for RQ1
%%   - "Robustness" list's "XGBoost hyperparameter sensitivity (RQ1 baseline)" bullet: outcome is DECIDED -- adopted into production, not just a sensitivity check
%%   - Limitations' "XGBoost's hyperparameter recipe is shared across research questions" paragraph + \fc{check its was tuned only to rq3}: now RQ1 is also individually tuned, only RQ2b/RQ3 share the fixed config
%% Search spec to add: n_estimators in {300,500,800}, max_depth in {3,4,6}, learning_rate in {0.02,0.04,0.08}, 5-fold spatial_block_kfold, selected on validation R2 (never test), per cohort
%% Architecture-size sensitivity (Robustness list bullet) outcome, for a one-line mention: null result -- no size helps both cohorts for any of the 3 guided-network variants; DNN somewhat size-sensitive, PINN/PINN_k 5-10x less so; val_loss ranking does not track test-R2 ranking (signature of noise, not a real capacity effect) -- left out of Results write-up as low-value, not reported there in detail.

%------------------------ rq3_ADDTIONS ------------------------
%% need to add tukey fence scan for rq3 outlier sorting
%% symptote fraction
%% Category-removal Moran's I check (Sec eval-diagnostics, "as a category-removal check for RQ3") outcome, for a one-line mention: stand_structure dominates category attribution, but removing terrain (the #2 category) LOWERS residual clustering while stand_structure removal doesn't -- the opposite of the working hypothesis, unresolved, no test distinguishes the competing explanations -- left out of Results write-up as an unresolved contradiction, not reported there in detail.

%------------------------ disclosed_limitations_ADDITIONS ------------------------
% DISCLOSED LIMITATION (add to main text/appendix): GNNWR's AIC/AICc are approximated via plain
% feature count rather than the true GWR-corrected effective degrees of freedom, and its F1/F2/F3
% significance tests are unavailable, because the package's exact hat-matrix computation would
% need ~54GB at this project's row counts; R2/RMSE (what GNNWR is actually compared against
% EN/XGBoost on throughout) are computed separately and are unaffected.
